\subsection{Ray tracing simulations}

\textbf{You will:}

\begin{itemize}
    \item define a cylindrical mirror for sagittal focusing
    \item define “externally” the optical element radius of curvature
    \item verify the best reflectivity is pbtained at demagnification 3
    \item include anticlastic bending
\end{itemize}

\textbf{Questions:}

\textbf{i)} Define a bending magnet source (e.q. ESRF1: $E_e$=6 GeV, $I$=0.2 A, $\sigma_x$=395$\mu$m, $\sigma_z$=36.9$\mu$m,$\sigma_x'$=9.87$\mu$rad,$\sigma_z'$=1.057$\mu$rad, horizontal acceptance 5 mrad). Set the photon energy interval in $E$=2000$\pm$10 eV. Define a (+,-) flat-flat Si111 double crystal monochromator tuned at 20 keV. Duplicate the second crystal and introduce the sagittal curvature to focus in the horizontal plane at the sample position. $p$=30~m, $q$=10~m. Calculate and compare the horizontal spot size. 

\textbf{ii)} With the help of a script, calculate the intensity vs magnification factor $M$. Fix $p$=30~m and set $q=Mp$. Verify that ratio $M=1/3$ is the optimum.

\textbf{Hints:} You may use the workspace files {\tt sagittalfocusing.ows} that also includes the python script that scans the magnification. One can see that the intensity peaks at about 1/3. A script is included to calculate the curvature radius: Rs (20 keV, M=1/3)=148.3 cm. 

iii) Study the anticlastic curvatura. For that, use the toroidal shape, with Major radius the anticlastic, and minor radius the sagittal. To express that $R_m$ has opposite cuvature than $R_s$, the torus pole must be located at the ``upper/inner=convex/concave" location. Try several anticlastic radius, starting from infinity (a large value) and going down to $R_m \approx$ 7.5 m where the diffracted intensity is almost zero. 

\textbf{Answer}


\begin{figure}[H]
\label{fig:sf-fig1}
\caption{Image at focal plane by the flat-flat system (left) and the flat-sagittal system (right)}
\includegraphics[width=0.48\textwidth]{sagittal-focusing/fig1a.png}
\includegraphics[width=0.48\textwidth]{sagittal-focusing/fig1b.png}
\end{figure}


\begin{figure}[H]
\label{fig:sf-fig2}
\caption{Intensity (in arbitrary units) versus magnification factor M for a point and monochromatic (E= 20 keV) source placed at 30 m from the sagittaly bent crystals. The beam divergences is 5 mrad. We clearly observe the maximum of the transmission at M=0.33, as predicted by the theory.}
\includegraphics[width=0.99\textwidth]{sagittal-focusing/fig2.png}
\end{figure}

